\documentclass[12pt,a4paper]{article}

\usepackage[margin=1.4cm]{geometry}
\usepackage{array}
\usepackage{enumitem}
\usepackage{multicol}

\setlist[itemize]{
	topsep=0pt,
	partopsep=0pt,
	parsep=0pt,
	itemsep=0pt
}

\newcolumntype{L}[1]{>{\raggedright\arraybackslash}p{#1}}
\newcommand{\checkbox}{\fbox{\rule{0pt}{0.8ex}\rule{0.8ex}{0pt}}}

\begin{document}
	
	\begin{center}
		{\Large\bfseries C3--C4: ESCUCHAR A LOS USUARIOS}
		
		\vspace{0.1cm}
		
		Emprendimiento e Innovación Grado 10
		
		\vspace{0.3cm}
		
		\begin{tabular}{L{0.12\textwidth}L{0.63\textwidth}}
			Grado: & \rule{2cm}{0.4pt} \\
			Fecha: & \rule{2cm}{0.4pt} \\
			Nombre 1: & \rule{\linewidth}{0.4pt} \\
			Nombre 2: & \rule{\linewidth}{0.4pt} \\
			Nombre 3: & \rule{\linewidth}{0.4pt} \\
		\end{tabular}
	\end{center}
	
	\section{Identificar a los usuarios}
	
	Revisen la observación realizada en la Clase 2. Identifiquen a las personas que experimentan o se ven afectadas por la situación.
	
	\begin{center}
		\small
		\begin{tabular}{|L{0.28\textwidth}|L{0.55\textwidth}|}
			\hline
			\textbf{Pregunta} & \textbf{Nuestra respuesta} \\
			\hline
			¿Qué situación estamos investigando? & \rule{0pt}{0.7cm} \\
			\hline
			¿Quiénes la experimentan o se ven afectados por ella? & \rule{0pt}{0.7cm} \\
			\hline
			¿Por qué deberíamos preguntarles a estas personas? & \rule{0pt}{0.7cm} \\
			\hline
		\end{tabular}
	\end{center}
	
	\section{Diseñar la encuesta}
	
	Creen una encuesta breve para conocer las experiencias, necesidades y dificultades de los usuarios. Las preguntas deben ser claras, neutrales y estar enfocadas en la situación.
	
	\begin{center}
		\small
		\begin{tabular}{|L{0.28\textwidth}|L{0.55\textwidth}|}
			\hline
			\textbf{Elemento} & \textbf{Nuestro plan de encuesta} \\
			\hline
			¿Quién responderá la encuesta? & \rule{0pt}{0.6cm} \\
			\hline
			¿Qué queremos aprender? & \rule{0pt}{0.7cm} \\
			\hline
			¿Cómo recopilaremos las respuestas? & \rule{0pt}{0.6cm} \\
			\hline
		\end{tabular}
	\end{center}
	
	\vspace{0.15cm}
	
	\textbf{Escriban sus preguntas:}
	
	\begin{center}
		\small
		\begin{tabular}{|L{0.10\textwidth}|L{0.73\textwidth}|}
			\hline
			\textbf{\#} & \textbf{Pregunta de la encuesta} \\
			\hline
			1 & \rule{0pt}{0.65cm} \\
			\hline
			2 & \rule{0pt}{0.65cm} \\
			\hline
			3 & \rule{0pt}{0.65cm} \\
			\hline
			4 & \rule{0pt}{0.65cm} \\
			\hline
			5 & \rule{0pt}{0.65cm} \\
			\hline
			6 & \rule{0pt}{0.65cm} \\
			\hline
			7 & \rule{0pt}{0.65cm} \\
			\hline
		\end{tabular}
	\end{center}
	
	\newpage
	
	\textbf{Buenas preguntas para una encuesta:}
	
	\begin{tabular}{@{}p{0.4\linewidth}@{\hspace{0.04\linewidth}}p{0.55\linewidth}@{}}
		\begin{itemize}[leftmargin=0.6cm]
			\item Preguntan sobre experiencias reales;
			\item Son neutrales y claras;
			\item No sugieren una respuesta;
		\end{itemize}
		&
		\begin{itemize}[leftmargin=0.6cm]
			\item Ayudan a identificar necesidades o dificultades;
			\item Evitan pedir a los usuarios que elijan una solución.
		\end{itemize}
	\end{tabular}
	
	\vspace{3mm}
	
	\textbf{Ejemplo:} ``¿Con qué frecuencia experimenta esta situación?''
	
	\textbf{Eviten:} ``¿Esta situación es difícil o frustrante para usted?''
	
	\section{Realizar la encuesta}
	
	Realicen la encuesta con los usuarios seleccionados. Registren sus respuestas con la mayor precisión posible.
	
	\begin{center}
		\small
		\begin{tabular}{|L{0.22\textwidth}|L{0.61\textwidth}|}
			\hline
			\textbf{Elemento} & \textbf{Evidencia de los usuarios} \\
			\hline
			Usuario / perfil & \rule{0pt}{0.5cm} \\
			\hline
			Experiencia relevante & \rule{0pt}{0.7cm} \\
			\hline
			Afirmación importante & \rule{0pt}{0.9cm} \\
			\hline
			Necesidad o frustración & \rule{0pt}{0.9cm} \\
			\hline
		\end{tabular}
	\end{center}
	
	\textbf{Importante:} Registren lo que los usuarios realmente dicen. No cambien sus respuestas para que coincidan con su idea original.
	
	\section{Comparar la evidencia}
	
	Comparen lo que dijeron los usuarios con lo que observaron en la Clase 2.
	%
	\begin{center}
		\small
		\begin{tabular}{|L{0.33\textwidth}|L{0.33\textwidth}|L{0.33\textwidth}|}
			\hline
			\textbf{Nuestra observación} & \textbf{Lo que dijeron los usuarios} & \textbf{Lo que aprendimos} \\
			\hline
			\rule{0pt}{1.3cm} & \rule{0pt}{1.3cm} & \rule{0pt}{1.3cm} \\
			\hline
			\rule{0pt}{1.3cm} & \rule{0pt}{1.3cm} & \rule{0pt}{1.3cm} \\
			\hline
			\rule{0pt}{1.3cm} & \rule{0pt}{1.3cm} & \rule{0pt}{1.3cm} \\
			\hline
		\end{tabular}
	\end{center}
	
	Busquen evidencia que confirme, amplíe o contradiga sus observaciones.
	
	\section{¿Necesidad o solución?}
	
	Una necesidad está respaldada por evidencia sobre lo que las personas experimentan, necesitan o valoran. Una solución es algo que ustedes proponen.
	
	\begin{center}
		\small
		\begin{tabular}{|L{0.40\textwidth}|L{0.43\textwidth}|}
			\hline
			\textbf{Necesidad basada en evidencia} & \textbf{Posible solución} \\
			\hline
			Los usuarios necesitan pasar menos tiempo esperando. & Crear un sistema de reserva para autoservicio. \\
			\hline
			Los usuarios tienen dificultades para encontrar información. & Crear un sitio web con función de búsqueda. \\
			\hline
		\end{tabular}
	\end{center}
	
	\begin{center}
		\small
		\begin{tabular}{|L{0.30\textwidth}|L{0.53\textwidth}|}
			\hline
			\textbf{Pregunta} & \textbf{Nuestro análisis} \\
			\hline
			¿Qué necesidad está respaldada por evidencia? & \rule{0pt}{1.0cm} \\
			\hline
			¿Qué solución estamos asumiendo? & \rule{0pt}{1.0cm} \\
			\hline
			¿Qué evidencia respalda la necesidad? & \rule{0pt}{1.0cm} \\
			\hline
		\end{tabular}
	\end{center}
	
	\vspace{-0.39cm}
	
	\section{Revisar nuestra comprensión}
	
	\vspace{-0.09cm}
	Utilicen la evidencia de las observaciones y de los usuarios para decidir si su comprensión original debe cambiar.
	
	\begin{center}
		\small
		\begin{tabular}{|L{0.27\textwidth}|L{0.27\textwidth}|L{0.27\textwidth}|}
			\hline
			\textbf{Idea original} & \textbf{Evidencia} & \textbf{¿Qué debería cambiar?} \\
			\hline
			\rule{0pt}{1.4cm} & \rule{0pt}{1.4cm} & \rule{0pt}{1.4cm} \\
			\hline
			\rule{0pt}{1.4cm} & \rule{0pt}{1.4cm} & \rule{0pt}{1.4cm} \\
			\hline
		\end{tabular}
	\end{center}
	
	\textbf{Nuestra conclusión:}
	
	\begin{itemize}[leftmargin=0.6cm]
		\item \checkbox\quad Nuestra comprensión original está respaldada.
		\item \checkbox\quad Nuestra comprensión necesita ser modificada.
		\item \checkbox\quad Nuestra comprensión ha cambiado sustancialmente.
		\item \checkbox\quad Necesitamos más información.
	\end{itemize}
	
	\vspace{0.25cm}
	
	\textbf{Evidencia más sólida:}
	
	\noindent\rule{\linewidth}{0.4pt}
	
	\vspace{0.25cm}
	
	\noindent\rule{\linewidth}{0.4pt}
	
	\vspace{0.25cm}
	
	\textbf{Necesidad más importante de los usuarios:}
	
	\noindent\rule{\linewidth}{0.4pt}
	
	\vspace{0.25cm}
	
	\noindent\rule{\linewidth}{0.4pt}
	
	\vspace{0.25cm}
	
	\textbf{Algo que aprendimos y que no esperábamos:}
	
	\noindent\rule{\linewidth}{0.4pt}
	
	\vspace{0.25cm}
	
	\noindent\rule{\linewidth}{0.4pt}
	
	\vspace{-0.29cm}
	
	\section{Verificación final}
	
	\begin{tabular}{@{}p{0.6\linewidth}@{\hspace{0.04\linewidth}}p{0.31\linewidth}@{}}
		\noindent\checkbox\quad Identificamos a los usuarios relevantes.
		
		\noindent\checkbox\quad Hicimos preguntas abiertas y registramos las respuestas reales.
		
		\vspace{0.12cm}
		
		\noindent\checkbox\quad Comparamos la evidencia de los usuarios con nuestras observaciones.
		
		\vspace{0.12cm}
		
		\noindent\checkbox\quad Identificamos una necesidad basada en evidencia.
		
		\vspace{0.12cm}
		
		\noindent\checkbox\quad Separamos las necesidades de las posibles soluciones.
		
		\vspace{0.12cm}
		
		\noindent\checkbox\quad Utilizamos evidencia para revisar nuestra comprensión original.
		&
		\begin{center}
			\textbf{Escuchen a los usuarios. Utilicen evidencia. Identifiquen la necesidad antes de proponer una solución.}
		\end{center}
	\end{tabular}
	
\end{document}
